% !TeX program = lualatex
% Compile twice: lualatex open_problems_500.tex
% All references and prose are embedded; no BibTeX database or external inputs.
\documentclass[11pt,a4paper]{article}
\usepackage[margin=24mm,headheight=14pt]{geometry}
\usepackage{fontspec}
\setmainfont{Latin Modern Roman}
\setsansfont{Latin Modern Sans}
\setmonofont{Latin Modern Mono}
\usepackage{amsmath,amssymb,mathtools}
\usepackage{unicode-math}
\setmathfont{Latin Modern Math}
\usepackage{microtype}
\usepackage{enumitem,xurl,needspace}
\usepackage[dvipsnames]{xcolor}
\usepackage{hyperref}
\hypersetup{unicode=true,colorlinks=true,linkcolor=MidnightBlue,urlcolor=MidnightBlue,
 pdftitle={Research notebook: Problems 40--49},pdfauthor={Source-annotated compilation},bookmarksnumbered=true}
\usepackage{bookmark}
\usepackage{tocloft}
\setlength{\cftsecnumwidth}{3em}
\cftsetpnumwidth{2.5em}
\cftsetrmarg{4em}
\usepackage{titlesec}
\titleformat{\section}{\Large\bfseries\raggedright}{\thesection}{0.7em}{}
\usepackage{fancyhdr}
\pagestyle{fancy}\fancyhf{}
\fancyhead[L]{\small\sffamily Mathematical problem statements}
\fancyhead[R]{\small Source edition 22}
\fancyfoot[C]{\thepage}
\setlength{\parindent}{0pt}
\setlength{\parskip}{5pt plus 1pt}
\setlength{\emergencystretch}{3em}
\setcounter{tocdepth}{1}
\setcounter{secnumdepth}{1}
\setlist[itemize]{leftmargin=3em,itemsep=5pt,topsep=4pt,parsep=0pt}
\allowdisplaybreaks
\newcommand{\N}{\mathbb{N}}
\newcommand{\Z}{\mathbb{Z}}
\newcommand{\Q}{\mathbb{Q}}
\newcommand{\R}{\mathbb{R}}
\newcommand{\C}{\mathbb{C}}
\newcommand{\F}{\mathbb{F}}
\newcommand{\A}{\mathbb{A}}
\newcommand{\PP}{\mathbb P}
\newcommand{\E}{\mathbb{E}}
\newcommand{\eps}{\varepsilon}
\newcommand{\dd}{\,\mathrm d}
\newcommand{\sref}[2]{\hyperlink{source:#1:#2}{[S#2]}}
\newcommand{\eref}[2]{\hyperlink{extra:#1:#2}{[E#2]}}
% Turn the next switch off for a shorter reading copy. The quotations preserve
% every exactTarget field, including qualifications omitted from a short summary.
\newif\ifshowcatalogue
\showcataloguetrue
\newcommand{\cataloguescope}[1]{\ifshowcatalogue
 \paragraph{Exact scope in the supplied catalogue.}
 \begin{quote}\small #1\end{quote}\fi}

% Research additions dated 27 September 2026; compile twice with LuaLaTeX.
\numberwithin{equation}{section}
\begin{document}
\setcounter{section}{46}
% ================================================================
% problemId: problem.the-fontaine-mazur-conjecture-on-geometric-galois-representations
% source releaseRank: 47; source releaseStatus: open_with_solved_subcases
% exactTarget SHA-256: 58e4a79bc122d5e31d6112e4dd17170c101ecf4a90d6cc252cbf3f99412cfbb4
\clearpage
\section{\texorpdfstring{The Fontaine-Mazur Conjecture on Geometric Galois Representations}{The Fontaine-Mazur Conjecture on Geometric Galois Representations}}\label{47}
\subsection{Definitions and mathematical statement}
Let $G_K=\operatorname{Gal}(\bar K/K)$ and $\rho:G_K\to\operatorname{GL}(V)$ be an irreducible continuous representation over a finite extension of $\Q_\ell$. Being unramified outside a finite set means inertia acts trivially at all remaining finite places. At $v\mid\ell$, being de Rham means the de Rham period module has the full dimension of $V$. The assertion is that
\[
 V\text{ is a subquotient of }H^i_{\mathrm{\acute et}}(X_{\bar K},\Q_\ell)\otimes\Q_\ell(m)
\]
with coefficients extended to those of $V$, for some smooth projective $X/K$, $i\ge0$, and $m\in\Z$. Arbitrary Tate twists are allowed; no nonnegative Hodge--Tate-weight condition is imposed.
\par\smallskip\textit{Formulation and concept references:} \sref{47}{1}.
\cataloguescope{Let K be a number field and l a prime. Every irreducible continuous finite-dimensional l-adic representation rho of Gal(Kbar/K), unramified outside finitely many places and de Rham at every place above l, occurs as a subquotient of H\textasciicircum{}i\_et(X\_Kbar)(m) for some smooth projective variety X over K, integer i >= 0, and integer Tate twist m. Etale cohomology is taken with the same l-adic coefficient field as rho. There is no nonnegative Hodge-Tate weight restriction.}
\subsection{Short English statement}
Must every irreducible Galois representation with the specified arithmetic and local geometric properties arise from algebraic geometry?
% BEGIN RESEARCH ADDITION 47
\subsection{Research attempt}

\paragraph{Current frontier.}
\textit{Literature and scope check, 27 September 2026.}
In \sref{47}{1}, printed page 13, Conjecture 3.2 and Remark 3.4 explain
exactly how removing the sign restriction on Hodge--Tate weights and allowing
Tate twists recovers the Fontaine--Mazur formulation. Those freedoms are
retained here. Established potential-automorphy theorems have additional
regularity, polarization and residual-image hypotheses; theorems for a
compatible system do not automatically apply to an arbitrary single
representation \eref{47}{1}.

The August 2026 manuscript \sref{47}{3}, Theorem B, treats two-dimensional
representations over totally real fields that are potentially crystalline
and ordinary with weights $\{0,1\}$ and totally odd determinant. It proves
potential modularity even in cases with reducible residual representation.
The text immediately after that theorem separately discusses an additional
local condition guaranteeing geometric realization. Theorem D proves
modularity over $\Q$ under the stated ordinary weight-$\{0,1\}$ hypotheses.
These are scoped results, not the all-fields, all-ranks, arbitrary-weight
target above. No independent audit of that recent manuscript is claimed.

\paragraph{Attack route.}
There are three distinct tasks: find a geometric realization after a finite
extension; bring it back to the original number field; and ensure that the
realization uses the precise smooth projective category in the question.
Potential automorphy addresses part of the first task under special
hypotheses, but calling a representation automorphic does not by itself
supply the required cohomology. We isolate and completely prove the latter
two tasks, including a connectedness safeguard. We also test whether twisting
or base extension removes the hypotheses that prevent applying the known
potential-automorphy results.

\paragraph{Attempt.}
\textbf{1. Descent from a genuinely geometric finite extension.}
Let $E/\Q_\ell$ be the coefficient field, $L/K$ a finite extension, and
$V$ a finite-dimensional continuous $E$-representation of $G_K$. No
irreducibility assumption is needed for the following implication:
\begin{equation}\label{eq:47:potential}
 \left.V\right|_{G_L}\text{ is a subquotient of }
 H^i_{\mathrm{\acute et}}(Y_{\bar L},E)(m)
 \quad\Longrightarrow\quad
 V\text{ has a smooth projective realization over }K,
\end{equation}
where $Y/L$ is smooth projective. The cohomological degree and twist on the
right may change if a geometrically connected variety is required.

Put $G=G_K$, $H=G_L$, and $d=[G:H]$. Finite-index induction is a direct sum
of $d$ copies of the underlying vector space, hence is exact. The tensor
identity gives
\[
 \operatorname{Ind}_H^G\operatorname{Res}_H^G V
       \simeq V\otimes_E E[G/H].
\]
The maps
\[
 v\longmapsto v\otimes\sum_{gH\in G/H}[gH],\qquad
 v\otimes\sum_{gH}a_{gH}[gH]\longmapsto
          \frac{1}{d}\Bigl(\sum_{gH}a_{gH}\Bigr)v
\]
are $G$-equivariant and their composite is the identity. Thus $V$ is a
direct summand of the induced restriction. Division by $d$ is legitimate in
$E$ even if $\ell\mid d$; this is not an integral-lattice assertion.

Regard the scheme $Y$ as a $K$-scheme $Z$ by composing its structural map
with $\operatorname{Spec}L\to\operatorname{Spec}K$. Smoothness and
projectivity are preserved because the latter map is finite etale.
After base change to $\bar K$, its components are the conjugates of $Y$,
and their cohomology, with the permutation of components and the transported
actions, is precisely
\begin{equation}\label{eq:47:ind-cohom}
 H^i_{\mathrm{\acute et}}(Z_{\bar K},E)(m)
 \simeq \operatorname{Ind}_{G_L}^{G_K}
       H^i_{\mathrm{\acute et}}(Y_{\bar L},E)(m).
\end{equation}
Finite disjoint-union cohomology and base change justify this identification;
no automorphic descent theorem is being invoked. These cohomological
functorialities are as in \eref{47}{2}.

If $\operatorname{Res}V=U/U'$ for $U'\subseteq U\subseteq W$, exact
induction exhibits $\operatorname{Ind}\operatorname{Res}V$ as
$\operatorname{Ind}U/\operatorname{Ind}U'$ inside a subquotient of
$\operatorname{Ind}W$. Taking the inverse image of the summand $V$ in this
quotient shows that $V$ itself is a subquotient of the right side of
\eqref{eq:47:ind-cohom}. This proves \eqref{eq:47:potential} in the smooth
projective scheme convention without any semisimplicity assertion for
arbitrary Galois cohomology.

\textbf{2. A geometrically connected realization without a category loophole.}
The preceding $Z$ need not be geometrically connected. To avoid resting the
argument on a permissive meaning of ``variety'', embed it into $\PP^N_K$
with constant codimension $c\ge2$, increasing $N$ if necessary. Here $Y$
may be taken pure-dimensional, as in a realization by a variety. Let
\[
 X=\operatorname{Bl}_{Z}(\PP^N_K),\quad
 j:D=\PP(N_{Z/\PP^N})\hookrightarrow X,\quad q:D\to Z,
 \quad \xi=c_1(\mathcal O_D(1)).
\]
The blow-up is smooth and projective; after extending scalars it remains
integral, since it is the blow-up of the integral projective space along a
proper smooth center. In particular, $X$ is geometrically connected.
The normal bundle of $D$ in $X$ is $\mathcal O_D(-1)$.

Using rational etale cohomology, define
\begin{align*}
 A &: H^i(Z_{\bar K},E)(m)\longrightarrow
          H^{i+2}(X_{\bar K},E)(m+1),& A&=j_*q^*,\\
 B &: H^{i+2}(X_{\bar K},E)(m+1)\longrightarrow
          H^i(Z_{\bar K},E)(m),& B&=-q_*(\xi^{c-2}j^*(-)).
\end{align*}
The Gysin self-intersection formula and projective-bundle integration give
\[
 BA=-q_*\bigl(\xi^{c-2}(-\xi)q^*(-)\bigr)
     =q_*(\xi^{c-1}q^*(-))=\operatorname{id}.
\]
The normalization is $q_*(\xi^{c-1})=1$ on every projective-space fiber.
These are the usual Gysin, projection and Chern-class identities
\eref{47}{2}. All maps are defined over $K$ and hence Galois-equivariant.
Thus the old cohomology is a direct summand of one cohomology group of a
geometrically connected smooth projective $X$, with the displayed changed
degree and twist. Both changes are allowed by the original target.
This supplies an explicit splitting, rather than assuming that cohomology
of a disconnected scheme is automatically an admissible answer.

\textbf{3. A complete, but already familiar, finite-image sector.}
If an irreducible representation $\tau$ has finite image, let $L/K$ be the
finite Galois extension cut out by its kernel, with group $\Gamma$.
The module $H^0((\operatorname{Spec}L)_{\bar K},E)$ is $E[\Gamma]$.
Every irreducible $E[\Gamma]$-module occurs in the regular representation:
it is a quotient, and averaging a linear splitting over the finite group
makes it a direct summand. Consequently, $\tau(m)$ is realized in
$H^0(\operatorname{Spec}L)(m)$ for every integer $m$. Construction 2 can
replace this finite etale scheme by a geometrically connected variety.
This proves the required conclusion whenever $V$ is a Tate twist of a
finite-image representation, without an oddness requirement. It is a known
special case, not a claim of new arithmetic progress.

\textbf{4. Why the known hypotheses cannot simply be normalized away.}
A Tate twist translates all Hodge--Tate numbers by the same integer
(with the sign determined by the convention). Differences and multiplicities
are unchanged. Thus repeated weights cannot become distinct by twisting.
After extending the local field, labeled weights at an embedding are
repeated at the embeddings above it, not arbitrarily replaced by new ones.
An appeal to regular potential automorphy still needs a regularity hypothesis.

There is a similar elementary obstruction to forcing self-duality by a
uniform twist. An essentially self-dual representation with a de Rham
multiplier must have its local weight multiset invariant under
$h\mapsto a-h$, where $a$ is the multiplier's weight. The multiset
$\{0,1,3\}$ has no such symmetry: pairing its extremes requires $a=3$,
while its middle weight requires $a=2$. Translation by an integer $t$
changes these requirements to $3+2t$ and $2+2t$, still inconsistent.
This is a test of a weight-based normalization lemma, not a construction of
a global irreducible representation violating Fontaine--Mazur. The
self-duality and distinct-weight conditions in \eref{47}{1} really are
additional restrictions.

\paragraph{Stress test.}
Restriction to $G_L$ need not preserve irreducibility. The proof of
\eqref{eq:47:potential} deliberately assumes a realization of the whole
restriction and never silently applies an irreducible-case theorem to it.
Induction is exact because the index is finite; no averaging over an
infinite profinite group has been used. Nor have we inferred that a general
subquotient splits. Only the explicitly constructed finite-permutation
summand and the explicitly split blow-up summand are split.

The remaining arithmetic hypothesis is substantial: for every $V$ in the
original statement, find some finite $L/K$ and an actual smooth projective
$Y/L$ with the left side of \eqref{eq:47:potential}. Local de Rham comparison
is a necessary property of cohomology, not a construction of such a $Y$.
Potential modularity must also be connected to an appropriate cohomological
realization; the qualification following Theorem B of \sref{47}{3} makes
this distinction particularly relevant. No unavailable universal
polarization, regularity, compatible-system, or residual-image theorem is
inserted into that gap.

\needspace{14\baselineskip}
\paragraph{Outcome.}
\textbf{Outcome: Conditional reduction.}
Potential geometric realization suffices for the exact base-field problem.
The descent, coefficient-field division, subquotient bookkeeping, and
geometric-connectedness repair are fully supplied above. Finite-image Tate
twists are handled directly. These structural arguments are not asserted
to be historically new.

What remains is potential \emph{geometric} realization for every irreducible,
finitely ramified, de Rham representation in the question. At the level of
universal existence this potential assertion is equivalent to the target,
since a realization over $K$ is already a potential realization. The
reduction isolates where an arithmetic advance is needed; it does not make
that advance or prove the conjecture.
% END RESEARCH ADDITION 47
\subsection{Sources}
\begin{itemize}\raggedright
\item[\textnormal{[S1]}]\hypertarget{source:47:1}{} Sug Woo Shin and Nicolas Templier, On fields of rationality for automorphic representations, author version dated 2 July 2014.
\url{https://math.berkeley.edu/~swshin/FieldRat.pdf}.
{\small\textit{Catalogue formulation source.}}
\item[\textnormal{[S2]}]\hypertarget{source:47:2}{} James Newton, Modularity of Galois Representations and Langlands Functoriality.
\url{https://link.springer.com/article/10.1007/s41745-022-00305-0}.
{\small\textit{Catalogue research/status reference; not a proof endorsement.}}
\item[\textnormal{[S3]}]\hypertarget{source:47:3}{} Jack A. Thorne, Towards the Fontaine–Mazur conjecture for GL(2).
\url{https://arxiv.org/html/2608.07186v1}.
{\small\textit{Catalogue research/status reference; not a proof endorsement.}}
% BEGIN RESEARCH SOURCES 47
\item[\textnormal{[E1]}]\hypertarget{extra:47:1}{} Thomas Barnet-Lamb,
Toby Gee, David Geraghty and Richard Taylor, \emph{Potential automorphy and
change of weight}, Annals of Mathematics (2) 179 (2014), 501--609;
arXiv:1010.2561, introduction, Theorems A and C, and Section 4.5.
\url{https://arxiv.org/abs/1010.2561}.
{\small\textit{Established potential-automorphy results with explicit
additional hypotheses; consulted 27 September 2026.}}
\item[\textnormal{[E2]}]\hypertarget{extra:47:2}{} James S. Milne,
\emph{Lectures on Etale Cohomology}, author-hosted course notes,
Sections 16 (purity and Gysin), 17 (proper base change), 22 (projection
formula) and 23 (Chern classes), especially Theorem 23.2.
\url{https://www.jmilne.org/math/CourseNotes/LEC.pdf}.
{\small\textit{Cohomological conventions and functorialities used in the
induction and blow-up calculations; consulted 27 September 2026.}}
% END RESEARCH SOURCES 47
\end{itemize}

\end{document}
